% \newpage
\section{Conclusions}

{\modified{}This work revisited} the conventional approach to prompt-based image editing within diffusion models. Contrary to established methods that leverage both negative and positive branches of Classifier-Free Guidance for editing ($\pin$ and $\pout$, respectively), {\modified{}the presented investigation reveals} that reconstruction is not necessary. {\modified{}The DDIM inversion process alone contains} enough information for effective editing, thereby questioning the need for manipulating $\pin$ and $\pout$ attention maps. {\modified{}The proposed approach simplifies} the editing process by applying guidance exclusively to the $\pin$ branch, which not only streamlines the method but also yields better results, as confirmed by extensive comparative analyses.


{\modified{}The primary contribution lies} in introducing and validating self-attention guidance as a superior mechanism for image editing tasks. Through quantitative analyses, ablation studies, and user feedback, {\modified{}it was demonstrated} that self-attention guidance captures a broader contextual understanding, enabling better edits compared to traditional cross-attention techniques. This approach preserves closer fidelity to the original image content while accurately applying the desired edits. 

{\modified{}A comprehensive comparative analysis}, supported by an extensive user study, shows that SAGE is preferred over every other compared method with statistical significance (\textit{p-value} $< 0.01$), achieving an average preference higher than 60.7\% among participants. This substantial margin underscores SAGE's effectiveness and potential to redefine standard practices in image editing with diffusion models.

% {\modified{}Several key areas} for improvement and future research {\modified{}have been identified}. For object removal, masking out self-attention guidance could prevent structure preservation in emptied areas, addressing issues in Sec.~\ref{sec:sup_limitations_structure} of the supplementary materials.
% {\modified{}Further exploration of sampling schedulers beyond DDIM and applying SAGE to distilled models \cite{meng2023distillation} and rectified-flow models \cite{liu2023flow} could benefit performance, taking advantage of more stable inference trajectoriesand higher quality models \cite{deng2024fireflow}}.
% Additionally, minimizing reliance on precise hyperparameter tuning by dynamically adjusting parameters based on task complexity, input characteristics, and denoising magnitude could enhance robustness across diverse editing scenarios.


{\modified{}Several directions for improvement and future research have been identified. For object removal, masking out self-attention guidance could mitigate unintended structure preservation. Investigating alternative sampling schedulers beyond DDIM and extending SAGE to higher-quality models, such as distilled models~\cite{meng2023distillation} or rectified-flow models~\cite{liu2023flow,deng2024fireflow}, may further enhance results while reducing the number of sampling steps required. Additionally, dynamically adapting hyperparameters based on task complexity, input characteristics, and denoising magnitude could improve robustness across a wide range of editing scenarios.}
